\documentclass{amsart}
% For dealing with references we use the comment environment
\usepackage{verbatim}
\newenvironment{reference}{\comment}{\endcomment}
%\newenvironment{reference}{}{}
\newenvironment{slogan}{\comment}{\endcomment}
\newenvironment{history}{\comment}{\endcomment}

% For commutative diagrams we use Xy-pic
\usepackage[all]{xy}

% We use 2cell for 2-commutative diagrams.
\xyoption{2cell}
\UseAllTwocells

% We use multicol for the list of chapters between chapters
\usepackage{multicol}

% This is generally recommended for better output
\usepackage{lmodern}
\usepackage[T1]{fontenc}

% For cross-file-references

% Package for hypertext links:
\usepackage{hyperref}

% For any local file, say "hello.tex" you want to link to please

% Theorem environments.
%
\theoremstyle{plain}
\newtheorem{theorem}[subsection]{Theorem}
\newtheorem{proposition}[subsection]{Proposition}
\newtheorem{lemma}[subsection]{Lemma}

\theoremstyle{definition}
\newtheorem{definition}[subsection]{Definition}
\newtheorem{example}[subsection]{Example}
\newtheorem{exercise}[subsection]{Exercise}
\newtheorem{situation}[subsection]{Situation}

\theoremstyle{remark}
\newtheorem{remark}[subsection]{Remark}
\newtheorem{remarks}[subsection]{Remarks}

\numberwithin{equation}{subsection}

% Macros
%
\def\lim{\mathop{\mathrm{lim}}\nolimits}
\def\colim{\mathop{\mathrm{colim}}\nolimits}
\def\Spec{\mathop{\mathrm{Spec}}}
\def\Hom{\mathop{\mathrm{Hom}}\nolimits}
\def\Ext{\mathop{\mathrm{Ext}}\nolimits}
\def\SheafHom{\mathop{\mathcal{H}\!\mathit{om}}\nolimits}
\def\SheafExt{\mathop{\mathcal{E}\!\mathit{xt}}\nolimits}
\def\Sch{\mathit{Sch}}
\def\Mor{\mathop{\mathrm{Mor}}\nolimits}
\def\Ob{\mathop{\mathrm{Ob}}\nolimits}
\def\Sh{\mathop{\mathit{Sh}}\nolimits}
\def\NL{\mathop{N\!L}\nolimits}
\def\CH{\mathop{\mathrm{CH}}\nolimits}
\def\proetale{{pro\text{-}\acute{e}tale}}
\def\etale{{\acute{e}tale}}
\def\QCoh{\mathit{QCoh}}
\def\Ker{\mathop{\mathrm{Ker}}}
\def\Im{\mathop{\mathrm{Im}}}
\def\Coker{\mathop{\mathrm{Coker}}}
\def\Coim{\mathop{\mathrm{Coim}}}

% Boxtimes
%
\DeclareMathSymbol{\boxtimes}{\mathbin}{AMSa}{"02}

%
% Macros for moduli stacks/spaces
%
\def\QCohstack{\mathcal{QC}\!\mathit{oh}}
\def\Cohstack{\mathcal{C}\!\mathit{oh}}
\def\Spacesstack{\mathcal{S}\!\mathit{paces}}
\def\Quotfunctor{\mathrm{Quot}}
\def\Hilbfunctor{\mathrm{Hilb}}
\def\Curvesstack{\mathcal{C}\!\mathit{urves}}
\def\Polarizedstack{\mathcal{P}\!\mathit{olarized}}
\def\Complexesstack{\mathcal{C}\!\mathit{omplexes}}
% \Pic is the operator that assigns to X its picard group, usage \Pic(X)
% \Picardstack_{X/B} denotes the Picard stack of X over B
% \Picardfunctor_{X/B} denotes the Picard functor of X over B
\def\Pic{\mathop{\mathrm{Pic}}\nolimits}
\def\Picardstack{\mathcal{P}\!\mathit{ic}}
\def\Picardfunctor{\mathrm{Pic}}
\def\Deformationcategory{\mathcal{D}\!\mathit{ef}}

\setlength{\emergencystretch}{3em}
\begin{document}
\title[Stacks Project, Tag 058G]{065 --- Lazard characterization of flat modules}
\maketitle
\noindent Copyright (C) 2005--2025 Johan de Jong. Stacks Project authors. Source: \href{https://stacks.math.columbia.edu/tag/058G}{Tag 058G}.

\noindent Editorial difficulty note (not author text). Difficulty H2: The proof turns finite relations in a flat module into factorizations through finite free modules. These factorizations construct a directed system whose colimit is the original module, proving the nontrivial direction of Lazard characterization..

\noindent Editorial provenance note (not author text). History: excerpt from revision \texttt{a0444\allowbreak 6e57e\allowbreak c1fbc\allowbreak 252a8\allowbreak 71afc\allowbreak ec775\allowbreak 2fb28\allowbreak 07b14}, algebra.tex, lines 19675--19740. Reference presentation was adapted; original-author context and core support blocks were retained or added. The mathematical wording is unchanged. Licensed under GNU FDL 1.2 or later; no invariant sections or cover texts. See ../licenses/COPYING. Context blocks reproduce author definitions and supporting results; they are not additional counted problems.

\begin{lemma}
\label{lemma-flat-factors-fp}
Let $M$ be an $R$-module.  Then $M$ is flat if and only if the following
condition holds: if $P$ is a finitely presented $R$-module and $f: P
\to M$ a module map, then there is a free finite $R$-module $F$ and
module maps $h: P \to F$ and $g: F \to M$ such that $f = g
\circ h$.
\end{lemma}

\begin{proof}
This is just a reformulation of condition (4) from
Lemma \ref{lemma-flat-factors-free}.
\end{proof}


\begin{lemma}
\label{lemma-flat-factors-free}
Let $M$ be an $R$-module.  The following are equivalent:
\begin{enumerate}
\item $M$ is flat.
\item If $f: R^n \to M$ is a module map and $x \in \Ker(f)$, then there
are module maps $h: R^n \to R^m$ and $g: R^m \to M$ such that
$f = g \circ h$ and $x \in \Ker(h)$.
\item Suppose $f: R^n \to M$ is a module map, $N \subset \Ker(f)$ any
submodule, and $h: R^n \to R^{m}$ a map such that $N \subset \Ker(h)$
and $f$ factors through $h$.  Then given any $x \in \Ker(f)$ we can find a map
$h': R^n \to R^{m'}$ such that $N + Rx \subset \Ker(h')$ and $f$
factors through $h'$.
\item If $f: R^n \to M$ is a module map and $N \subset \Ker(f)$ is a
finitely generated submodule, then there are module maps $h: R^n \to
R^m$ and $g: R^m \to M$ such that $f = g \circ h$ and $N \subset
\Ker(h)$.
\end{enumerate}
\end{lemma}

\begin{proof}
That (1) is equivalent to (2) is just a reformulation of the equational
criterion for flatness\footnote{In fact, a module map $f : R^n \to M$
corresponds to a choice of elements $x_1, x_2, \ldots, x_n$ of $M$
(namely, the images of the standard basis elements $e_1, e_2, \ldots,
e_n$); furthermore, an element $x \in \Ker(f)$ corresponds to a
relation between these $x_1, x_2, \ldots, x_n$ (namely, the relation
$\sum_i f_i x_i = 0$, where the $f_i$ are the coordinates of $x$).
The module map $h$ (represented as an $m \times n$-matrix)
corresponds to the matrix $(a_{ij})$ from Lemma \ref{lemma-flat-eq},
and the $y_j$ of Lemma \ref{lemma-flat-eq} are the images of the
standard basis vectors of $R^m$ under $g$.}.   To show
(2) implies (3), let $g: R^m \to M$
be the map such that $f$ factors as $f = g \circ h$.  By (2) find $h'': R^m
\to R^{m'}$ such that $h''$ kills $h(x)$ and $g: R^m \to M$
factors through $h''$.  Then taking $h' = h'' \circ h$ works.  (3) implies (4)
by induction on the number of generators of $N \subset \Ker(f)$ in (4).
Clearly (4) implies (2).
\end{proof}


\begin{lemma}[Equational criterion of flatness]
\label{lemma-flat-eq}
A module $M$ over $R$ is flat if and only if
every relation in $M$ is trivial.
\end{lemma}

\begin{proof}
Assume $M$ is flat and let $\sum f_i x_i = 0$ be a relation in $M$.
Let $I = (f_1, \ldots, f_n)$, and let
$K = \Ker(R^n \to I, (a_1, \ldots, a_n) \mapsto \sum_i a_i f_i)$.
So we have the short exact sequence
$0 \to K \to R^n \to I \to 0$. Then $\sum f_i \otimes x_i$
is an element of $I \otimes_R M$ which maps
to zero in $R \otimes_R M = M$. By flatness
$\sum f_i \otimes x_i$ is zero in $I \otimes_R M$.
Thus there exists an element of $K \otimes_R M$ mapping
to $\sum e_i \otimes x_i \in R^n \otimes_R M$ where $e_i$
is the $i$th basis element of $R^n$.
Write this element as $\sum k_j \otimes y_j$
and then write the image of $k_j$ in $R^n$ as
$\sum a_{ij} e_i$ to get the result.

\medskip\noindent
Assume every relation is trivial, let $I$
be a finitely generated ideal, and let $x = \sum f_i \otimes x_i$
be an element of $I \otimes_R M$ mapping to zero in $R \otimes_R M = M$.
This just means exactly that $\sum f_i x_i$ is a relation in
$M$. And the fact that it is trivial implies easily that
$x$ is zero, because
$$
x
=
\sum f_i \otimes x_i
=
\sum f_i \otimes \left(\sum a_{ij}y_j\right)
=
\sum \left(\sum f_i a_{ij}\right) \otimes y_j
=
0
$$
\end{proof}


\begin{theorem}[Lazard's theorem]
\label{theorem-lazard}
Let $M$ be an $R$-module.  Then $M$ is flat if and only if it is the colimit of
a directed system of free finite $R$-modules.
\end{theorem}

\begin{proof}
A colimit of a directed system of flat modules is flat, as taking directed
colimits is exact and commutes with tensor product. Hence if $M$ is the colimit
of a directed system of free finite modules then $M$ is flat.

\medskip\noindent
For the converse, first recall that any module $M$ can be written as the
colimit of a directed system of finitely presented modules, in the following
way.  Choose a surjection $f: R^I \to M$ for some set $I$, and let $K$
be the kernel.  Let $E$ be the set of ordered pairs $(J, N)$ where $J$ is a
finite subset of $I$ and $N$ is a finitely generated submodule of $R^J \cap K$.
Then $E$ is made into a directed partially ordered set by defining $(J, N) \leq
(J', N')$ if and only if $J \subset J'$ and $N \subset N'$.  Define $M_e =
R^J/N$ for $e = (J, N)$, and define $f_{ee'}: M_e \to M_{e'}$ to be
the natural map for $e \leq e'$.  Then $(M_e, f_{ee'})$ is a directed system and
the natural maps $f_e: M_e \to M$ induce an isomorphism
$\colim_{e
\in E} M_e \xrightarrow{\cong} M$.

\medskip\noindent
Now suppose $M$ is flat.  Let $I = M \times \mathbf{Z}$, write $(x_i)$ for
the canonical basis of $R^{I}$, and take in the above discussion $f: R^I
\to M$ to be the map sending $x_i$ to the projection of $i$ onto $M$.
To prove the theorem it suffices to show that the $e \in E$ such that $M_e$
is free form a cofinal subset of $E$.  So let $e = (J, N) \in E$ be arbitrary.
By Lemma \ref{lemma-flat-factors-fp} there is a free finite module $F$ and maps
$h: R^J/N \to F$ and $g: F \to M$ such that the natural map
$f_e: R^J/N \to M$ factors as $R^J/N \xrightarrow{h} F
\xrightarrow{g} M$.  We are going to realize $F$ as $M_{e'}$ for some $e' \geq
e$.

\medskip\noindent
Let $\{ b_1, \ldots, b_n \}$ be a finite basis of $F$.  Choose $n$ distinct
elements $i_1, \ldots, i_n \in I$ such that $i_{\ell} \notin J$ for all $\ell$,
and such that the image of $x_{i_{\ell}}$ under $f: R^I \to M$ equals
the image of $b_{\ell}$ under $g: F \to M$.  This is possible since
every element of $M$ can be written as $f(x_i)$ for infinitely many
distinct $i \in I$ (by our choice of $I$).  Now let
$J' = J \cup \{i_1, \ldots , i_n \}$, and define $R^{J'}
\to F$ by $x_i \mapsto h(x_i)$ for $i \in J$ and $x_{i_{\ell}} \mapsto
b_{\ell}$ for $\ell = 1, \ldots, n$.  Let $N' = \Ker(R^{J'} \to F)$.
Observe:
\begin{enumerate}
\item The square
$$
\xymatrix{
R^{J'} \ar[r] \ar@{^{(}->}[d] & F \ar[d]^{g} \\
R^{I} \ar[r]_{f} & M
}
$$
is commutative,
%$R^{J'} \to F$ factors $f: R^I \to M$,
hence $N' \subset K = \Ker(f)$;
\item $R^{J'} \to F$ is a surjection onto a free finite module, hence
it splits and so $N'$ is finitely generated;
\item $J \subset J'$ and $N \subset N'$.
\end{enumerate}
By (1) and (2) $e' = (J', N')$ is in $E$, by (3) $e' \geq e$, and by
construction $M_{e'} = R^{J'}/N' \cong F$ is free.
\end{proof}
\end{document}
